% ============================================================
% PRÉAMBULE — Rapport ARGOS — Thème bleu professionnel
% ============================================================

% ---- Langue et encodage ----
\usepackage[utf8]{inputenc}
\usepackage[T1]{fontenc}
\usepackage[french]{babel}
\addto\captionsfrench{\renewcommand{\listfigurename}{Liste des figures}}
\usepackage{pdfpages}
\frenchbsetup{StandardLists=true}

% ---- Géométrie de page ----
\usepackage[a4paper,left=2.6cm,right=2.3cm,top=3cm,bottom=2.35cm,
            headheight=24pt,headsep=0.6cm]{geometry}
\usepackage[edges]{forest}  % Arborescences / graphes

% ---- Polices & Interligne ----
\usepackage{newtxmath}
\usepackage[scaled=1.04]{newtxtext} % Grossit la police de 3% pour matcher Word
\usepackage{helvet}             % Helvetica (Titre/sans-serif)
\usepackage{microtype}          % Micro-typographie (évite les débordements de texte)
\usepackage{setspace}
\setstretch{1.5}                % Mise en page plus compacte et lisible

% ---- Réglages de justification & prévention des débordements ----
\tolerance=1000
\emergencystretch=3em
\hfuzz=0.5pt

% ---- Unités, chiffres et graphiques ----
\usepackage{pgfplots}
\pgfplotsset{compat=1.18}
\usepackage{siunitx}
\usepackage{textcomp}
\usepackage{shorttoc}

% ---- Couleurs : thème bleu unique ----
\usepackage{xcolor}
\definecolor{argosnavy}{HTML}{102A43}     % bleu marine — texte fort
\definecolor{argosblue}{HTML}{1A365D}     % bleu principal — titres
\definecolor{argosmid}{HTML}{2C5282}      % bleu moyen — liens, accents
\definecolor{argoslight}{HTML}{B7D3F2}    % bleu clair soutenu — fonds de lignes
\definecolor{argospale}{HTML}{F3F7FC}     % bleu très pâle
\definecolor{argosrule}{HTML}{7EA6CF}     % bleu gris — filets
\definecolor{argosgray}{HTML}{4A5568}     % gris-bleu — texte secondaire
\definecolor{argostablehead}{RGB}{32,62,96}   % bleu nuit légèrement éclairci
\definecolor{argostablerow}{RGB}{183,211,242}
\definecolor{argostablerule}{RGB}{55,78,108}  % traits foncés et bien visibles


% ---- Tableaux ----
\usepackage{array}
\usepackage{tabularx}
\usepackage{longtable}
\usepackage{multirow}
\usepackage{colortbl}
\usepackage{xltabular}

\usepackage{booktabs}
\usepackage{etoolbox}
% Style commun à tous les tableaux du rapport.
\renewcommand{\arraystretch}{1.32}
\setlength{\extrarowheight}{1.2pt}
\setlength{\heavyrulewidth}{0.9pt}
\setlength{\lightrulewidth}{0.55pt}
\setlength{\cmidrulewidth}{0.45pt}
\setlength{\aboverulesep}{0pt}
\setlength{\belowrulesep}{0pt}
\renewcommand{\topfraction}{0.85}       % Autorise les figures à occuper 85% du haut de la page
\renewcommand{\bottomfraction}{0.75}    % Autorise les figures à occuper 75% du bas de la page
\renewcommand{\textfraction}{0.1}       % Laisse seulement 10% minimum pour le texte sur une page-figure
\renewcommand{\floatpagefraction}{0.8}  % Évite qu'une page ne soit réservée qu'aux figures si elle ne remplit pas 80%
\setlength{\tabcolsep}{4pt}
\setlength{\arrayrulewidth}{0.9pt}
\setlength{\LTleft}{0pt}
\setlength{\LTright}{0pt}
\arrayrulecolor{argostablerule}

\newcommand{\tableheadrow}{\rowcolor{argostablehead}}
\newcommand{\tablespecrow}{\rowcolor{argostablerow}}

% Grille et bandes de lecture uniformes : les lignes de données restent
% lisibles sur toute leur largeur, y compris dans les tableaux multipages.
\AtBeginEnvironment{tabular}{\setstretch{1}\arrayrulecolor{argostablerule}\rowcolors{2}{argostablerow}{white}\small\setlength{\tabcolsep}{4pt}}
\AtBeginEnvironment{tabularx}{\setstretch{1}\arrayrulecolor{argostablerule}\rowcolors{2}{argostablerow}{white}\small\setlength{\tabcolsep}{4pt}}
\AtBeginEnvironment{longtable}{\setstretch{1}\arrayrulecolor{argostablerule}\rowcolors{2}{argostablerow}{white}\footnotesize\setlength{\tabcolsep}{3.5pt}\setlength{\LTleft}{0pt}\setlength{\LTright}{0pt}\sloppy}
\AtBeginEnvironment{xltabular}{\setstretch{1}\arrayrulecolor{argostablerule}\rowcolors{2}{argostablerow}{white}\footnotesize\setlength{\tabcolsep}{3.5pt}\setlength{\LTleft}{0pt}\setlength{\LTright}{0pt}\sloppy}

\newcommand{\thc}[1]{\textcolor{white}{\textbf{#1}}}
\newcommand{\tablebottomrule}{\noalign{\hrule height 1.2pt}} 
\newcommand{\tablesourcepromoteurs}{%
  \par\smallskip{\raggedleft\footnotesize\itshape Source : Promoteurs{2025-2026}\par}}
\makeatletter
\newcommand{\tablerole}{%
  \par\smallskip{\raggedright\footnotesize\itshape
  \textbf{Rôle du tableau :} ce tableau constitue un support de synthèse pour
  l'étude de « \@currentlabelname ». Il organise les éléments nécessaires à la
  justification des choix, à la vérification des résultats et à la prise de
  décision dans la suite du rapport.\par}}
\makeatother

\usepackage{adjustbox}
\usepackage{pgf-pie}

% ---- Légendes des figures / tableaux ----
\usepackage{caption}
\usepackage{subcaption}
\captionsetup{
  font={small},
  labelfont={bf,color=argosblue},
  textfont={color=black},
  margin=10pt,
  belowskip=2pt,
  aboveskip=3pt
}
\captionsetup[table]{position=top}
\captionsetup[figure]{position=bottom}

% ---- Listes ----
\usepackage{enumitem}
\setlist[itemize,1]{label=\textcolor{argosblue}{\bfseries\textgreater}}
% Définition du nouveau style de puce flèche
\newcommand{\argospuce}{\textcolor{argosblue}{\textbf{$\gtr5$}}}

% ---- Images, pages paysage, TikZ ----
\usepackage{graphicx}
\graphicspath{{figures/}{figures/outils/}{figures/images/}}
\usepackage{pdflscape}
\usepackage{float}
\usepackage[section]{placeins}
\usepackage{flafter}
\usepackage{tikz}
\usetikzlibrary{arrows.meta,calc,fit,positioning,shapes.geometric,shapes.multipart}
\setlength{\textfloatsep}{10pt plus 2pt minus 3pt}
\setlength{\intextsep}{8pt plus 2pt minus 2pt}
\setlength{\floatsep}{7pt plus 2pt minus 2pt}

% Maintient chaque illustration à proximité immédiate de son point d'appel.
% Les barrières empêchent aussi plusieurs figures différées de se regrouper
% plus loin dans le document ou de franchir un titre de section.
\BeforeBeginEnvironment{figure}{\FloatBarrier}
\AfterEndEnvironment{figure}{\FloatBarrier}

% ---- Maths (Correction du conflit newtxmath / amssymb) ----
\usepackage{amsmath}
\let\Bbbk\undefined % Annule la définition préalable pour éviter l'erreur de compilation
\usepackage{amssymb}

% ---- Pages auxiliaires ----
\usepackage{needspace}
\usepackage{csquotes}

% ============================================================
% TITRES (HIÉRARCHIE VISUELLE SOBRE ET DISTINCTE)
% ============================================================
\usepackage{titlesec}

% Partie : Très grand (~29pt)
\titleformat{\part}[display]
  {\normalfont\sffamily\Huge\bfseries\color{argosblue}\centering}
  {\textcolor{argosmid}{\Large PARTIE \thepart}}
  {8pt}
  {\titlerule[1.2pt]\vspace{10pt}\MakeUppercase}

% Chapitre : compact et placé près du haut de la page
\titleformat{\chapter}[display]
  {\normalfont\sffamily\LARGE\bfseries\color{argosblue}\centering}
  {\textcolor{argosmid}{\normalsize CHAPITRE \thechapter}}
  {4pt}
  {\titlerule[1pt]\vspace{4pt}\MakeUppercase}
\titlespacing*{\chapter}{0pt}{-24pt}{14pt}

% Titres non numérotés sans barre décorative : résumé, abstract, sigles,
% introduction générale, conclusion et références.
\titleformat{name=\chapter,numberless}[display]
  {\normalfont\sffamily\LARGE\bfseries\color{argosblue}\centering}
  {}
  {0pt}
  {\MakeUppercase}

% Section : Moyen-Grand (~17pt) -> Modifié avec \MakeUppercase ici
\titleformat{\section}
  {\normalfont\sffamily\Large\bfseries\color{argosblue}\centering}
  {\thesection}{10pt}{\MakeUppercase}
\titlespacing*{\section}{0pt}{14pt plus 3pt minus 2pt}{8pt plus 2pt minus 1pt}

% Sous-section : Moyen (~14pt)
\titleformat{\subsection}
  {\normalfont\sffamily\large\bfseries\color{argosmid}}
  {\thesubsection}{8pt}{}
\titlespacing*{\subsection}{0pt}{11pt plus 2pt minus 2pt}{5pt plus 1pt minus 1pt}

% Sous-sous-section : Gras (~12pt)
\titleformat{\subsubsection}
  {\normalfont\sffamily\normalsize\bfseries\color{argosnavy}}
  {\thesubsubsection}{8pt}{}
\titlespacing*{\subsubsection}{0pt}{8pt plus 2pt minus 1pt}{3pt plus 1pt minus 1pt}

% ============================================================
% En-têtes et pieds de page
% ============================================================
\usepackage{fancyhdr}
\newcommand{\currentpartmark}{}
\newcommand{\argossetpart}[1]{%
  \gdef\currentpartmark{#1}%
  \markboth{\currentpartmark}{}%
}
\newcommand{\argossetchapter}[1]{%
  \markboth{\currentpartmark}{#1}%
}
\renewcommand{\chaptermark}[1]{%
  \argossetchapter{#1}%
}
\renewcommand{\sectionmark}[1]{}

% Pages précédant la liste des sigles : pied de page uniquement.
\fancypagestyle{argosfront}{%
  \fancyhf{}
  \fancyfoot[R]{\small\sffamily\color{argosblue}\thepage}
  \renewcommand{\headrulewidth}{0pt}
  \renewcommand{\footrulewidth}{0pt}
}
\fancypagestyle{plain}{%
  \fancyhf{}
  \fancyfoot[R]{\small\sffamily\color{argosblue}\thepage}
  \renewcommand{\headrulewidth}{0pt}
  \renewcommand{\footrulewidth}{0pt}
}

% En-tête et pied de page principaux du rapport.
\newcommand{\argosgradientrule}{%
  \hbox{%
    \color{argospale}\rule{0.15\headwidth}{1.2pt}%
    \color{argoslight}\rule{0.20\headwidth}{1.2pt}%
    \color{argosnavy}\rule{0.30\headwidth}{1.2pt}%
    \color{argoslight}\rule{0.20\headwidth}{1.2pt}%
    \color{argospale}\rule{0.15\headwidth}{1.2pt}%
  }%
}
% La mise en forme est appliquée ici, et non dans les marques LaTeX :
% les titres restent ainsi visibles après les changements de page, les
% flottants et les pages en orientation paysage.
\newcommand{\argosheadertext}[1]{%
  \sffamily\footnotesize\bfseries\color{argosnavy}\nouppercase{#1}%
}
\fancypagestyle{argosmain}{%
  \fancyhf{}
  \fancyhead[L]{\argosheadertext{\leftmark}}
  \fancyhead[R]{\argosheadertext{\rightmark}}
  \fancyfoot[L]{\sffamily\small\bfseries\color{argosnavy}ARGOS}
  \fancyfoot[R]{\small\sffamily\color{argosblue}\thepage}
  \renewcommand{\headrulewidth}{1.2pt}
  \renewcommand{\footrulewidth}{1.2pt}
  \renewcommand{\headrule}{\argosgradientrule\vskip-1.2pt}
  \renewcommand{\footrule}{\argosgradientrule\vskip2pt}
}
\newcommand{\argosactivateheaders}{%
  \pagestyle{argosmain}%
  \fancypagestyle{plain}{%
    \fancyhf{}
    \fancyhead[L]{\argosheadertext{\leftmark}}
    \fancyhead[R]{\argosheadertext{\rightmark}}
    \fancyfoot[L]{\sffamily\small\bfseries\color{argosnavy}ARGOS}
    \fancyfoot[R]{\small\sffamily\color{argosblue}\thepage}
    \renewcommand{\headrulewidth}{1.2pt}
    \renewcommand{\footrulewidth}{1.2pt}
    \renewcommand{\headrule}{\argosgradientrule\vskip-1.2pt}
    \renewcommand{\footrule}{\argosgradientrule\vskip2pt}%
  }%
}

% Catalogue unique : chaque entrée alimente à la fois la bibliographie finale
% et les notes de bas de page. Les URL enregistrées dans le catalogue ne sont
% jamais imprimées : les appels \cite affichent la référence complète en note.
\makeatletter
\newcommand{\DeclareArgosReference}[3]{%
  \expandafter\gdef\csname argosbib@title@#2\endcsname{#3}%
  \expandafter\gdef\csname argosbib@label@#2\endcsname{#1}%
}
\newcommand{\DeclareArgosUrl}[2]{%
  \expandafter\gdef\csname argosbib@url@#1\endcsname{#2}%
}
\newcommand{\argosbibitem}[1]{%
  \item[\hbox{[\csname argosbib@label@#1\endcsname]}]%
  \begingroup\let\url\@gobble
    \csname argosbib@title@#1\endcsname
  \endgroup
}
\newcommand{\argosonefootcite}[1]{%
  \ifcsname argosbib@title@#1\endcsname
    \footnote{\begingroup\let\url\@gobble
      \csname argosbib@title@#1\endcsname
    \endgroup}%
  \else
    \PackageWarning{argos}{Référence bibliographique inconnue : #1}%
  \fi
}
\newcommand{\bibfootcite}[1]{\forcsvlist{\argosonefootcite}{#1}}
\renewcommand{\cite}[1]{\bibfootcite{#1}}
% Catalogue documentaire unique du rapport ARGOS.
% Format : \DeclareArgosReference{indice}{clé}{référence complète avec URL éventuelle}

% URL seules utilisées dans les notes de bas de page.
\DeclareArgosUrl{unctad2025}{https://unctad.org/publication/review-maritime-transport-2025}
\DeclareArgosUrl{cts2025}{https://containerstatistics.com/december-2025-press-release/}
\DeclareArgosUrl{gmi2025}{https://www.gminsights.com/industry-analysis/shipping-containers-market}
\DeclareArgosUrl{omi2026}{https://www.imo.org/en/mediacentre/pressbriefings/pages/raft-of-shipping-rules-in-force-from-1-january-2026.aspx}
\DeclareArgosUrl{pad2025}{https://www.diaspocamtv.com/news/details/transport-maritime-7-372-mouvements-de-navires-enregistres-dans-les-ports-de-douala-et-kribi-entre-2021-et-2025}
\DeclareArgosUrl{corridor2025}{https://lefisk.cm/blog/transit-marchandises-cameroun-tchad-rca-ttu-corridors}
\DeclareArgosUrl{pak2026}{https://pak.cm/fr/port-de-kribi-la-performance-au-service-de-la-nation/}
\DeclareArgosUrl{imo_web}{https://www.imo.org/}
\DeclareArgosUrl{banquemondiale2025}{https://www.banquemondiale.org/fr/topic/trade/publication/the-african-continental-free-trade-area}
\DeclareArgosUrl{tresor2025}{https://www.tresor.economie.gouv.fr/Pays/CM/indicateurs-et-conjoncture}
\DeclareArgosUrl{cemac2025}{https://cemac.int/wp-content/uploads/2025/12/Rapport-de-Surveillance-Multilaterale-2024-et-Perspectives-2025-et-2026.pdf}
\DeclareArgosUrl{agenceecofin}{https://www.agenceecofin.com/actualites-infrastructures/2204-137774-infrastructures-routieres-1-6-milliard-pour-rehabiliter-la-section-camerounaise-du-corridor-douala-bangui}
\DeclareArgosUrl{flask}{https://flask.palletsprojects.com/}
\DeclareArgosUrl{corridor_transit}{https://lefisk.cm/blog/transit-marchandises-cameroun-tchad-rca-ttu-corridors}
\DeclareArgosUrl{postgresql}{https://www.postgresql.org/docs/16/}
\DeclareArgosUrl{playwright}{https://playwright.dev/python/docs/intro}
\DeclareArgosUrl{tailwind}{https://tailwindcss.com/docs}
\DeclareArgosUrl{esp32}{https://www.espressif.com/en/products/socs/esp32-s3}
\DeclareArgosUrl{ublox}{https://www.u-blox.com/en/product/neo-m8u-module}
\DeclareArgosUrl{sim7000}{https://www.simcom.com/}
\DeclareArgosUrl{lcsc}{https://www.lcsc.com/}
\DeclareArgosUrl{leaflet}{https://leafletjs.com/}
\DeclareArgosUrl{chartjs}{https://www.chartjs.org/}
\DeclareArgosUrl{mermaid}{https://mermaid.js.org/}
\DeclareArgosUrl{logcluster}{https://lca.logcluster.org/23-central-african-republic-reseau-routier}
\DeclareArgosUrl{art}{http://www.art.cm/}
\DeclareArgosUrl{antic}{http://www.antic.cm/}
\DeclareArgosUrl{bpifrance}{https://bpifrance-creation.fr/encyclopedie/structures-juridiques/choix-du-statut-generalites/apports-societe}
\DeclareArgosUrl{vernimmen_dalloz}{https://www.librairiedalloz.fr/livre/9782247208562-finance-d-entreprise-edition-2022-pascal-quiry-yann-le-fur-pierre-vernimmen/}
\DeclareArgosUrl{vernimmen_fnac}{https://www.fnac.com/a15983411/Yann-Le-Fur-Finance-d-entreprise-2022-20ed}
\DeclareArgosUrl{iso21502}{https://committee.iso.org/sites/tc258/home/projects/published/iso-21502.html}
\DeclareArgosUrl{cncc_missions}{https://www.cncc.cm/fr/a-propos/historique}
\DeclareArgosUrl{apn_missions}{https://www.apn.cm/missions-et-vision-de-apn/}
\DeclareArgosUrl{ohada_auscgie}{https://www.ohada.org/droit-des-societes-commerciales-et-du-gie/}
\DeclareArgosUrl{cgi_cameroun}{https://www.impots.cm/fr/actualites/code-general-des-impots-mis-jour-au-1er-janvier-2023}
\DeclareArgosUrl{cnps_employeurs}{https://www.cnps.cm/fr/employeurs/regles-generales-pour-les-employeurs1.html}
\DeclareArgosUrl{loi2024_donnees}{https://www.prc.cm/files/2b/9f/21/1055fa3c2251b4c4248fd301f584daaf.pdf}
\DeclareArgosUrl{minpmeesa_cfce}{https://www.minpmeesa.cm/site/inhoud/uploads/2023/11/sopecam2022.pdf}
\DeclareArgosUrl{art_homologation}{https://www.art.cm/fr/node/17334}
\DeclareArgosUrl{imo_solas}{https://www.imo.org/en/about/conventions/pages/international-convention-for-the-safety-of-life-at-sea-(solas),-1974.aspx}
\DeclareArgosUrl{imo_imdg}{https://www.imo.org/en/ourwork/safety/pages/dangerousgoods-default.aspx}
\DeclareArgosUrl{rgpd_officiel}{https://eur-lex.europa.eu/eli/reg/2016/679/oj?locale=fr}
\DeclareArgosUrl{incitations2025}{https://www.prc.cm/fr/actualites/actes/lois/8095-loi-n-2025-015-du-17-decembre-2025-portant-ratification-de-l-ordonnance-n-2025-002-du-18-juillet-2025-fixant-les-incitations-a-l-investissement-en-republique-du-cameroun}

\DeclareArgosReference{1}{unctad2025}{UNCTAD, 2025. \textit{Review of Maritime Transport 2025}. Genève~: Nations Unies. \url{https://unctad.org/publication/review-maritime-transport-2025} (consulté le 2~juin~2026)}
\DeclareArgosReference{2}{cts2025}{Container Trades Statistics (CTS), 2025. \textit{December 2025 Press Release -- Volume des échanges conteneurisés mondiaux}. \url{https://containerstatistics.com/december-2025-press-release/} (consulté le 3~juin~2026)}
\DeclareArgosReference{3}{gmi2025}{Global Market Insights, 2025. \textit{Shipping Containers Market Size, Growth, Industry Analysis, 2025–2035}. \url{https://www.gminsights.com/industry-analysis/shipping-containers-market} (consulté le 4~juin~2026)}
\DeclareArgosReference{4}{omi2026}{Organisation Maritime Internationale (OMI), 2026. \textit{Résolutions MSC.550(108) et MEPC.384(81) -- Amendements aux conventions internationales SOLAS (règle V/31-32) et MARPOL (Protocole I art. V) relatifs au signalement obligatoire des pertes de conteneurs en mer}, en vigueur au 1\textsuperscript{er}~janvier~2026. \url{https://www.imo.org/en/mediacentre/pressbriefings/pages/raft-of-shipping-rules-in-force-from-1-january-2026.aspx} (consulté le 5~juin~2026)}
\DeclareArgosReference{5}{pad2025}{PAKAZURE, 2026. \textit{Analyse consolidée du trafic maritime camerounais entre 2021 et 2025}, relayée par Diaspocam TV le 1\textsuperscript{er}~février~2026. \url{https://www.diaspocamtv.com/news/details/transport-maritime-7-372-mouvements-de-navires-enregistres-dans-les-ports-de-douala-et-kribi-entre-2021-et-2025} (consulté le 28~mai~2026)}
\DeclareArgosReference{6}{corridor2025}{Plateforme Logistique des Corridors, 2025. \textit{Transit des marchandises sur le corridor Douala-N'Djamena-Bangui et dépendance des pays enclavés}. \url{https://lefisk.cm/blog/transit-marchandises-cameroun-tchad-rca-ttu-corridors} (consulté le 6~juin~2026)}
\DeclareArgosReference{7}{pak2026}{Port Autonome de Kribi (PAK), 2026. \textit{Port de Kribi : la performance au service de la Nation}, publié le 21~avril~2026. \url{https://pak.cm/fr/port-de-kribi-la-performance-au-service-de-la-nation/} (consulté le 24~juin~2026)}
\DeclareArgosReference{20}{iso2018}{Organisation Internationale de Normalisation (ISO), 2018. \textit{Norme ISO~18185:2018 -- Conteneurs du commerce -- Sceaux électroniques (E-Seals) -- Parties~1 à~5}. Genève~: Secrétariat central de l'ISO (consulté le 8~septembre~2026).}
\DeclareArgosReference{22}{imo_web}{International Maritime Organization (IMO)~: conventions de sécurité maritime SOLAS et réglementation du fret connecté. \url{https://www.imo.org/} (consulté le 7~juin~2026)}
\DeclareArgosReference{23}{wsc2025}{World Shipping Council (WSC), 2025. \textit{Report on Containers Lost at Sea -- 2025 Update} (données statistiques et enjeux de la traçabilité active du fret). Washington D.C.~: WSC Communications (consulté le 8~septembre~2026).}
\DeclareArgosReference{8}{banquemondiale2025}{Banque mondiale, 2025. \textit{The African Continental Free Trade Area: Economic and Distributional Effects}. \url{https://www.banquemondiale.org/fr/topic/trade/publication/the-african-continental-free-trade-area} (consulté le 8~juin~2026)}
\DeclareArgosReference{9}{tresor2025}{Direction générale du Trésor (France) et BEAC, 2025. \textit{Cameroun -- Indicateurs et conjoncture 2025} et \textit{Rapport annuel 2024}. \url{https://www.tresor.economie.gouv.fr/Pays/CM/indicateurs-et-conjoncture} (consulté le 9~juin~2026)}
\DeclareArgosReference{10}{cemac2025}{Commission de la CEMAC, 2025. \textit{Rapport de Surveillance Multilatérale 2024 et Perspectives 2025-2026}. \url{https://cemac.int/wp-content/uploads/2025/12/Rapport-de-Surveillance-Multilaterale-2024-et-Perspectives-2025-et-2026.pdf} (consulté le 10~juin~2026)}
\DeclareArgosReference{11}{cncc2025}{Conseil National des Chargeurs du Cameroun (CNCC), 2025. \textit{Rapport annuel sur l'état des litiges et de la sécurité des marchandises sur les corridors de transit Douala-N'Djamena et Kribi-Bangui}. Douala~: Publications du CNCC (consulté le 8~septembre~2026).}
\DeclareArgosReference{12}{agenceecofin}{Agence Ecofin, 2025. \textit{Infrastructures et corridors -- Le Tchad se tourne vers la Banque mondiale pour couvrir 500 zones blanches et réhabiliter le corridor}. \url{https://www.agenceecofin.com/actualites-infrastructures/2204-137774-infrastructures-routieres-1-6-milliard-pour-rehabiliter-la-section-camerounaise-du-corridor-douala-bangui} (consulté le 11~juin~2026)}
\DeclareArgosReference{44}{supptic2024}{SUP'PTIC, 2024. \textit{Conception et réalisation d'une plateforme d'alerte en cas de détresse et de suivi des enfants (Helpscape~X)}. Mémoire de fin d'études d'ingénieur des travaux des télécommunications. Yaoundé~: Bibliothèque de SUP'PTIC (consulté le 8~septembre~2026).}
\DeclareArgosReference{19}{loi2010_tel}{République du Cameroun, 2010. \textit{Loi n°\,2010/013 du 21~décembre~2010 régissant les communications électroniques au Cameroun}. Yaoundé~: Journal Officiel (consulté le 8~septembre~2026).}
\DeclareArgosReference{25}{loi2010_cyber}{République du Cameroun, 2010. \textit{Loi n°\,2010/012 du 21~décembre~2010 relative à la cybersécurité et à la cybercriminalité au Cameroun}. Yaoundé~: Journal Officiel (consulté le 8~septembre~2026).}
\DeclareArgosReference{50}{babu_mbuangi2023}{F. BABU et A. MBUANGI-MAVUNGU, \textit{Manuel simplifié du droit des transports routier et maritime (OHADA et CEMAC)}, Le Lys Bleu Éditions, 2023 (consulté le 8~septembre~2026).}
\DeclareArgosReference{39}{flask}{Pallets Projects (Flask), \textit{Official Documentation for Python/Flask Framework}. \url{https://flask.palletsprojects.com/} (consulté le 12~juin~2026)}
\DeclareArgosReference{43}{corridor_transit}{Plateforme Logistique des Corridors, 2025. \textit{Transit des marchandises sur le corridor Douala-N'Djamena-Bangui et dépendance des pays enclavés}. \url{https://lefisk.cm/blog/transit-marchandises-cameroun-tchad-rca-ttu-corridors} (consulté le 6~juin~2026)}
\DeclareArgosReference{40}{postgresql}{PostgreSQL Global Development Group, \textit{PostgreSQL 16 Documentation \& pgcrypto Extension Architecture}. \url{https://www.postgresql.org/docs/16/} (consulté le 13~juin~2026)}
\DeclareArgosReference{45}{playwright}{Playwright for Python, \textit{End-to-End Visual UI Testing and Automation Framework}. \url{https://playwright.dev/python/docs/intro} (consulté le 14~juin~2026)}
\DeclareArgosReference{46}{tailwind}{Tailwind Labs, \textit{Tailwind CSS Utility-First Framework Documentation}. \url{https://tailwindcss.com/docs} (consulté le 15~juin~2026)}
\DeclareArgosReference{13}{esp32}{Espressif Systems, \textit{ESP32-S3 Series Microcontroller Hardware Design and Technical Reference}. \url{https://www.espressif.com/en/products/socs/esp32-s3} (consulté le 16~juin~2026)}
\DeclareArgosReference{15}{ublox}{u-blox AG, \textit{NEO-M8U Untethered Dead Reckoning (UDR) Module Technical Specifications}. \url{https://www.u-blox.com/en/product/neo-m8u-module} (consulté le 17~juin~2026)}
\DeclareArgosReference{14}{sim7000}{SIMCom Wireless Solutions, \textit{SIM7000G Cellular LTE-M and NB-IoT Module Datasheet \& MQTT Application Notes}. \url{https://www.simcom.com/} (consulté le 18~juin~2026)}
\DeclareArgosReference{47}{lcsc}{LCSC Electronics, catalogues d'approvisionnement des composants microélectroniques industriels. \url{https://www.lcsc.com/} (consulté le 19~juin~2026)}
\DeclareArgosReference{41}{leaflet}{Leaflet.js, \textit{Open-Source JavaScript Library for Mobile-Friendly Interactive Maps}. \url{https://leafletjs.com/} (consulté le 20~juin~2026)}
\DeclareArgosReference{42}{chartjs}{Chart.js, \textit{Simple yet Flexible JavaScript Charting for Designers \& Developers}. \url{https://www.chartjs.org/} (consulté le 21~juin~2026)}
\DeclareArgosReference{48}{mermaid}{Mermaid.js, \textit{Diagramming and Charting Tool for Documenting Architecture Code}. \url{https://mermaid.js.org/} (consulté le 22~juin~2026)}
\DeclareArgosReference{24}{logcluster}{Logistics Cluster, 2025. \textit{Digital Logistics Capacity Assessments -- Fiche RCA : Réseau routier et infrastructures}. \url{https://lca.logcluster.org/23-central-african-republic-reseau-routier} (consulté le 23~juin~2026)}
\DeclareArgosReference{21}{art}{Agence de Régulation des Télécommunications (ART), procédures d'homologation des équipements terminaux radio au Cameroun. \url{http://www.art.cm/} (consulté le 25~juin~2026)}
\DeclareArgosReference{49}{antic}{Agence Nationale des Technologies de l'Information et de la Communication (ANTIC), normes de sécurité et d'audit des plateformes Cloud applicatives. \url{http://www.antic.cm/} (consulté le 26~juin~2026)}
\DeclareArgosReference{36}{bpifrance}{Bpifrance Création, \textit{Choix du statut -- généralités : apports en société}. \url{https://bpifrance-creation.fr/encyclopedie/structures-juridiques/choix-du-statut-generalites/apports-societe} (consulté le 27~juin~2026)}
\DeclareArgosReference{37}{vernimmen_dalloz}{P. Vernimmen, P. Quiry et Y. Le Fur, \textit{Finance d'entreprise}, édition 2022, Dalloz. \url{https://www.librairiedalloz.fr/livre/9782247208562-finance-d-entreprise-edition-2022-pascal-quiry-yann-le-fur-pierre-vernimmen/} (consulté le 28~juin~2026)}
\DeclareArgosReference{38}{vernimmen_fnac}{Y. Le Fur, P. Quiry et P. Vernimmen, \textit{Finance d'entreprise 2022}, fiche bibliographique FNAC. \url{https://www.fnac.com/a15983411/Yann-Le-Fur-Finance-d-entreprise-2022-20ed} (consulté le 28~juin~2026)}
\DeclareArgosReference{16}{iso21502}{Organisation internationale de normalisation, \textit{ISO~21502:2020 -- Management de projets, programmes et portefeuilles : recommandations sur le management de projets}. \url{https://committee.iso.org/sites/tc258/home/projects/published/iso-21502.html} (consulté le 29~juin~2026)}
\DeclareArgosReference{17}{cncc_missions}{Conseil National des Chargeurs du Cameroun, \textit{Historique, vision et mission globale du CNCC}. \url{https://www.cncc.cm/fr/a-propos/historique} (consulté le 1\textsuperscript{er}~juin~2026)}
\DeclareArgosReference{18}{apn_missions}{Autorité Portuaire Nationale, \textit{Missions et vision de l'APN}. \url{https://www.apn.cm/missions-et-vision-de-apn/} (consulté le 2~juin~2026)}
\DeclareArgosReference{30}{ohada_auscgie}{OHADA, \textit{Acte uniforme relatif au droit des sociétés commerciales et du groupement d'intérêt économique}, adopté le 30~janvier~2014. \url{https://www.ohada.org/droit-des-societes-commerciales-et-du-gie/} (consulté le 3~juin~2026)}
\DeclareArgosReference{31}{cgi_cameroun}{Direction générale des impôts du Cameroun, \textit{Code général des impôts}, version officielle publiée par la DGI. \url{https://www.impots.cm/fr/actualites/code-general-des-impots-mis-jour-au-1er-janvier-2023} (consulté le 4~juin~2026)}
\DeclareArgosReference{33}{cnps_employeurs}{Caisse Nationale de Prévoyance Sociale du Cameroun, \textit{Règles générales pour les employeurs}. \url{https://www.cnps.cm/fr/employeurs/regles-generales-pour-les-employeurs1.html} (consulté le 5~juin~2026)}
\DeclareArgosReference{34}{loi2024_donnees}{République du Cameroun, \textit{Loi n°2024/017 du 23 décembre 2024 relative à la protection des données à caractère personnel}. \url{https://www.prc.cm/files/2b/9f/21/1055fa3c2251b4c4248fd301f584daaf.pdf} (consulté le 6~juin~2026)}
\DeclareArgosReference{35}{minpmeesa_cfce}{MINPMEESA, \textit{La création d'entreprise : formalités auprès des Centres de Formalités de Création d'Entreprises}. \url{https://www.minpmeesa.cm/site/inhoud/uploads/2023/11/sopecam2022.pdf} (consulté le 7~juin~2026)}
\DeclareArgosReference{29}{art_homologation}{Agence de Régulation des Télécommunications, \textit{Demande d'homologation d'un équipement de communications électroniques ou d'une installation radioélectrique}. \url{https://www.art.cm/fr/node/17334} (consulté le 8~juin~2026)}
\DeclareArgosReference{26}{imo_solas}{Organisation maritime internationale, \textit{Convention internationale de 1974 pour la sauvegarde de la vie humaine en mer (SOLAS)}. \url{https://www.imo.org/en/about/conventions/pages/international-convention-for-the-safety-of-life-at-sea-(solas),-1974.aspx} (consulté le 9~juin~2026)}
\DeclareArgosReference{27}{imo_imdg}{Organisation maritime internationale, \textit{Code maritime international des marchandises dangereuses (Code IMDG)}. \url{https://www.imo.org/en/ourwork/safety/pages/dangerousgoods-default.aspx} (consulté le 10~juin~2026)}
\DeclareArgosReference{28}{rgpd_officiel}{Union européenne, \textit{Règlement (UE) 2016/679 du 27 avril 2016 (RGPD)}. \url{https://eur-lex.europa.eu/eli/reg/2016/679/oj?locale=fr} (consulté le 11~juin~2026)}
\DeclareArgosReference{32}{incitations2025}{République du Cameroun, \textit{Loi n°2025/015 du 17 décembre 2025 portant ratification de l'ordonnance n°2025/002 fixant les incitations à l'investissement}. \url{https://www.prc.cm/fr/actualites/actes/lois/8095-loi-n-2025-015-du-17-decembre-2025-portant-ratification-de-l-ordonnance-n-2025-002-du-18-juillet-2025-fixant-les-incitations-a-l-investissement-en-republique-du-cameroun} (consulté le 12~juin~2026)}

\renewenvironment{thebibliography}[1]{%
  \list{}% 
    {\setlength{\leftmargin}{3em}%
     \setlength{\itemindent}{0pt}%
     \setlength{\labelwidth}{2.5em}%
     \setlength{\labelsep}{0.7em}%
     \setlength{\itemsep}{0.5em}%
     \setlength{\parsep}{0pt}%
     \setlength{\topsep}{0pt}%
     \setlength{\partopsep}{0pt}%
     \setlength{\listparindent}{0pt}%
     \setlength{\rightmargin}{0pt}}%
  \sloppy
  \clubpenalty4000
  \widowpenalty4000
  \sfcode`\.\@m}
  {\endlist}
\makeatother

% ============================================================
% Liste des figures / tableaux / Table des matières
% ============================================================
\usepackage{tocloft}
\renewcommand{\cftpartfont}{\sffamily\bfseries\color{argosblue}}
\renewcommand{\cftchapfont}{\sffamily\bfseries\color{argosblue}}
\renewcommand{\cftsecfont}{\sffamily\color{argosnavy}}
\renewcommand{\cftsubsecfont}{\sffamily\color{argosgray}}
\renewcommand{\cftchappagefont}{\sffamily\color{argosblue}}
\renewcommand{\cftpartpagefont}{\sffamily\color{argosblue}}
\renewcommand{\cftfigfont}{\sffamily\color{argosnavy}}
\renewcommand{\cfttabfont}{\sffamily\color{argosnavy}}
\renewcommand{\cftfigpagefont}{\sffamily\color{argosgray}}
\renewcommand{\cfttabpagefont}{\sffamily\color{argosgray}}
\setlength{\cftbeforechapskip}{6pt}

% ---- AJOUT POUR AFFICHER LES POINTS DE SUITE (Sommaire/Table des matières) ----
\renewcommand{\cftpartleader}{\color{argosblue}\cftdotfill{\cftdotsep}}
\renewcommand{\cftchapleader}{\color{argosblue}\cftdotfill{\cftdotsep}}

% ============================================================
% Hyperliens & References
% ============================================================
% Les tableaux xltabular sont mesurés en plusieurs passes. Des noms d'ancrage
% fondés uniquement sur les compteurs produiraient alors des destinations PDF
% dupliquées (table.4.7, table.4.8, etc.). Les identifiants séquentiels évitent
% ces collisions tout en conservant les liens, renvois et signets visibles.
\usepackage[colorlinks=true,linkcolor=argosmid,citecolor=argosmid,urlcolor=argosmid,
            bookmarksopen=true,bookmarksnumbered=true,pdfborder={0 0 0},
            hypertexnames=false]{hyperref}
\usepackage{xurl}
\usepackage[capitalise,french]{cleveref}

% ============================================================
% Encadrés sobres
% ============================================================
\usepackage{mdframed}

\mdfdefinestyle{argosbox}{%
  linecolor=argosmid,
  linewidth=1.1pt,
  leftline=true,rightline=false,topline=false,bottomline=false,
  backgroundcolor=argospale,
  innertopmargin=8pt,innerbottommargin=8pt,
  innerleftmargin=11pt,innerrightmargin=11pt,
  skipabove=11pt,skipbelow=11pt
}
\newenvironment{encadre}[1]
 {\begin{mdframed}[style=argosbox]\sffamily\bfseries\color{argosblue}#1\par\nopagebreak\smallskip\normalfont\rmfamily\color{black}}
 {\end{mdframed}}

\newcommand{\code}[1]{\texttt{\small #1}}

% ============================================================
% Numérotation et profondeur
% ============================================================
\counterwithin{figure}{chapter}
\counterwithin{table}{chapter}

\setcounter{tocdepth}{2}
\setcounter{secnumdepth}{3}
